\documentclass[10pt]{article}

\usepackage[utf8]{inputenc}

\usepackage[T1]{fontenc}

\usepackage{amsmath}

\usepackage{amsfonts}

\usepackage{amssymb}

\usepackage[version=4]{mhchem}

\usepackage{stmaryrd}

\usepackage{bbold}

\usepackage{graphicx}

\usepackage[export]{adjustbox}

\graphicspath{ {./images/} }



\begin{document}

П а лю тин Е. А., Мателатическая логика М., 1979; [5] Handbook of mathematical logic, Amst., 1977; [6] D e v l i n K. J., Aspects of constructibility, B.-[a.o.],1973. B. Н. гришин.



СКОЛЬЖЕНИИ ГРУППА р е г у л я р н о г о н ак р ы т и я $p: X \rightarrow Y-$ группа $\Gamma(p)$ таких гомеоморфизмов $\vee$ пространства $X$ на себя, что $\gamma p=p(X$ и $Y-$ связные локально линейно связные хаусдорфовы топологич. Пространства). Так, С. г. накрытия окружности действительной прямой $\mathbb{R}$, определяемого формулой $t \mapsto(\cos t, \sin t)$, будет групіа сдвигов $t \mapsto t+2 \pi n, n \in \mathbb{Z}$.



Группа $\Gamma(p)$ является дискретной аруппой преобразований пространства $X$, действующей свободно (т. е. $\gamma(x)=x \Rightarrow \gamma=1)$, и пространство $Y$ естественно изоморфно факторпространству $X / \Gamma(p)$. С. г. $\Gamma(p)$ изоморфна факторгруппе бундаментальной группь $\pi_{1}\left(Y, y_{0}\right)$, где $y_{0} \in Y$, по образу группы $\pi_{1}\left(X, x_{0}\right)$, где $p\left(x_{0}\right)=y_{0}$, при гомоморфизме, индуцированном отображением $p$. В частности, для универсального накрытия $p$ C. г. пзоморфна фундаментальной группс пространства $Y$. Лum.: X у С ы-Щ з я н, Теория гомотопий, пер. с англ., М.,

1964.



СІОЛЬЗЯІЕГО СРЕДНЕГО ПРОЦЕСС - стационарный в широком смысле случайный процесс, к-рый может быть получен с помощью применения нек-рого линейного преобразования к процессу с некоррелированными значениями (т. е. к процессу белого шума). Часто С. с. І1. наз. также более частный процесс $X(t)$ с дискретным временем $t=0, \pm 1, \cdots$, представимый в виде



\[

X(t)=Y(t)+b_{1} Y(t-1)+\ldots+b_{q} Y(t-q)

\]



где $\boldsymbol{E I}^{Y}(t)=0, \mathrm{E} Y(t) \boldsymbol{Y}(s)=\sigma^{2} \delta_{t s}, \delta_{t s}$-символ Кронекера (так что $Y(t)$-процесс белого пІума со спектральной плотностью $\left.\sigma^{2} / 2 \pi\right), q$-нек-рое целое положительное число, а $b_{1}, \ldots, b_{q}$ - постоянные коәффициенты. Спектральная плотность $f(\lambda)$ такого С. с. П. определяется формулой



\[

\begin{gathered}

f(\lambda)=\left(\sigma^{2} / 2 \pi\right)\left|\psi\left(e^{i \lambda}\right)\right|^{2}, \\

\psi(z)=b_{0}+b_{1} z+\ldots+b_{q^{z}} q, b_{0}=1,

\end{gathered}

\]



a его коррельяцинная функция $r(k)=E X(t) X(t-k)$ имеет вид



\[

\begin{gathered}

r(k)=\sigma^{2} \sum_{j=0}^{q-|k|} b_{j} b_{j+\mid} k \quad \text { при }|k| \leqslant q, \\

r(k)=0 \text { при }|k|>q .

\end{gathered}

\]



Обратно, если корреляционная функция $r(k)$ стационарного процесса $X(t)$ с дискретным временем $t$ обладает тем свойством, что $r(k)=0$ при $|k|>q$ для какого-то целого положительного $q$, то $X(t)$ - әто С. с. п. порядка $q$, т. е. он допускает представление вида (1), где $\boldsymbol{Y}(t)$ - белый шум (см., напр., [1], $\$ .7)$.



Наряду со С. с. п. конечного порядка $q$, представимыми в виде (1), существуют также два типа С. с. п. с дискретным временем бесконечного порядка, а именно: односторонние С. с. п., допускающие представление вида



\[

X(t)=\sum_{j=0}^{\infty} b_{j} Y(t-j)

\]



где $Y(t)$ - белый шум, а ряд в правой части (2) сходится в среднем квадратичном (и, значит, $\sum_{j=0}^{\infty}\left|b_{j}\right|^{2}<$ $<\infty)$, и более общие двусторонние С. с. п., представимые в виде



\[

X(t)=\sum_{j=-\infty}^{\infty} b_{j} Y(t-j),

\]



где $Y(t)$ - белый пгум, а $\sum_{j=-\infty}^{\infty}\left|b_{j}\right|^{2}<\infty$. Класс двусторонних С. с. п. совпадает с классом стационарных процессов $X(t)$, пмеющих спектральную плотность $f(\lambda)$, а класс односторонних С. с. п.- с классом про- цессов, имеющих спектральную плотность $f(\lambda)$ такую, чTo



\[

\int_{-\pi}^{\pi} \log f(\lambda) d \lambda>-\infty

\]



(см. [2], [1], [3]).



Односторонним или соответственно двусторонним C. с. п. с непрерывным временем наз. стационарный процесс $X(t),-\infty<t<\infty$, представимый в виде



\[

X(t)=\int_{0}^{\infty} b(s) d Y(t-s), \int_{0}^{\infty}|b(s)|^{2} d s<\infty

\]



или соответственно в виде



\[

X(t)=\int_{-\infty}^{\infty} b(s) d Y(t-s), \int_{-\infty}^{\infty}|b(s)|^{2} d s<\infty

\]



\includegraphics[max width=\textwidth, center]{2023_07_08_c2394bcea4e929bdcb72g-1}

белого шума. Класс двусторонних С. с. п. с непрерывным временем совпадает с классом стационарных процессов $X(t)$, имеющих спектральную плотноств $f(\lambda)$, а класс односторонних С. с. п. с непрерывным временем-с классом процессов, имеющих такую спектральную плотность $f(\lambda)$, что



\[

\int_{-\infty}^{\infty} \log f(\lambda)\left(1+\lambda^{2}\right)^{-1} d \lambda>-\infty

\]



(см. [4], [3], [5]).



Jит.: [1] А н д е р с о н Т., Статистический анализ временных рядов, пер. с англ., М., 1976; [2] К о л м о г о р о в А. Н.,

"Бюлл. Моск. гос. ун-та", 1941, т. 2, в. 6, с. 1 А 0 ; [3] Д. у б Дж., Вероятностные процессы, пер. с англ., M., 1956; [4] K a r h u-

n e n K., "Ann. Acad. Sci. Fennicae. Ser. A. Math.-Phys.", 1947, n e n K., "Ann. Acad. Sci. Fennicae. Ser. A. Math.-Phys.", 1947,

$\mathcal{N}_{0}$ 37, p. 3-79; [5] p о з а н о в Ю. А., Стационарные случайные процессы, М., 1963 . А. М. Яглом.



СКОЛЬЗЯІИЙ ВЕКТОР - см. Ве ктор.



СКРЕПІЕННОЕ ПРОИЗВЕДЕНИЕ Г р У П П Ы $G$ И к о л ц а $K$ - ассоциативное кольцо, определяемое следующей конструкцией. Пусть заданы однозначноө отображение $\sigma$ группы $G$ в группу автоморфизмов ассоциативного кольца $K$ с единицей и семейство



\[

\rho=\left\{\rho_{g, h} \mid g, h \in G\right\}

\]



обратимых элементов кольца $K$, удовлетворяющее условиям



\[

\begin{gathered}

\rho_{g_{1}, g_{2}, g_{3}} \cdot \rho_{g_{2}, g_{3}}=\rho_{g_{1} g_{2}, g_{3}} \rho_{g_{1}, g_{2}}^{g_{3} \sigma} ; \\

\alpha^{g_{1} \sigma \cdot g_{2} \sigma}=\rho_{g_{1}, g_{2}}^{-1} \alpha^{\left(g_{1} g_{2}\right) \sigma} \rho_{g_{1}, g_{2}}

\end{gathered}

\]



для всех $\alpha \in K$ и $g_{1}, g_{2}, g_{3} \in G$. Семейство $\rho$ наз. с и с т емо й ф а к т о р о в. Элементами С. п. группы $G$ и кольца $K$ при системе факторов $\rho$ и отображении $\sigma$ будут всевозможные формальные конечные суммы вида



\[

\sum_{g \in G} \alpha_{g} t_{g}, g \in G, \alpha_{g} \in K

\]



$\left(t_{g}\right.$ - символ, однозначно сопоставляемый каждому элементу $g \in G)$, а операции определяются формулами



\[

\begin{gathered}

\sum_{g \in G} \alpha_{g} t_{g}+\sum_{g \in G} \beta_{g} t_{g}=\sum_{g \in G}\left(\alpha_{g}+\beta_{g}\right) t_{g} ; \\

\left(\sum_{g \in G} \alpha_{g} t_{g}\right)\left(\sum_{g \in G} \beta_{g} t_{g}\right)= \\

=\sum_{h \in G}\left(\sum_{x y=h, x, y \in G}\left(\alpha_{x} \beta_{y}^{x \sigma} \rho_{x, y}\right)\right) t_{h} .

\end{gathered}

\]



Это кольцо обозначается $K(G, \rho, \sigma)$; элементы $t_{g}$ образуют $K$-базис С. п. $K(G, \rho, \sigma)$.



Если $\sigma$ отображает $G$ в единичный автоморфизм кольца $K$, то С. п. $K(G, \rho)$ наз. с к р ещ е нвы м г р у пІІ о в ы м ко л ь цо м, а если, кроме того, $\rho_{g, h}=1$ для всех $g, h \in G$, то $K(G, \rho, \sigma)$ - г р у п по в о е ко л ьц о группы $G$ над кольцом $K$ (см. Групповая алгебра). Пусть $K$ - поле и $\sigma$ - мономорфизм. Тогда С. п. $K(G, \rho, \sigma)$ - простое кольцо, являющееся С. п. поля с его группой Галуа.



Jum.: [1] S e h g a l S. K., Topics in group rings, N. Y., 1978; [2] Б о в д и А. А., “Сиб. матем. ж.», 1963, т. /, с. 481 500 ; [3] Итоги науки и техники. Современные проблемы матема тики, т. 2, М., 1973 , с. $5-118$.



A. A. Бовди.





\end{document}